%
% ---------- header -----------------------------------------------------------
%
% project       kaneton
%
% license       kaneton
%
% file          /home/mycure/kaneton/view/curriculum/EPITA/epita.tex
%
% created       julien quintard   [sat oct 20 09:31:28 2007]
% updated       julien quintard   [sat dec  8 03:00:19 2007]
%

%
% ---------- path -------------------------------------------------------------
%

\def\path{../..}

%
% ---------- packages ---------------------------------------------------------
%

%
% template
%

%
% ---------- header -----------------------------------------------------------
%
% project       kaneton
%
% license       kaneton
%
% file          /home/mycure/kaneton/view/template/paper.tex
%
% created       julien quintard   [wed may 16 18:17:37 2007]
% updated       julien quintard   [fri oct  5 07:00:45 2007]
%

%
% class
%

\documentclass[10pt,a4wide]{article}

%
% packages
%

\usepackage[english]{babel}
\usepackage[T1]{fontenc}
\usepackage{a4wide}
\usepackage{fancyheadings}
\usepackage{multicol}
\usepackage{indentfirst}
\usepackage{graphicx}
\usepackage{color}
\usepackage{xcolor}
\usepackage{verbatim}
\usepackage{aeguill}

\pagestyle{fancy}

\setlength{\footrulewidth}{0.3pt}
\setlength{\parindent}{0.3cm}
\setlength{\parskip}{2ex plus 0.5ex minus 0.2ex}

%
% logos
%

\newcommand{\logos}
  {
    \begin{center}
      \includegraphics[scale=0.8]{\path/logo/kaneton.pdf}
    \end{center}
  }

%
% prototype
%

\newcommand\prototype[2]{
  \begin{tabular}{p{0.2cm}p{13.8cm}}
  & #1
  \end{tabular}

  \begin{tabular}{p{1cm}p{13cm}}
  & #2
  \end{tabular}}

%
% verbatim stuff
%

\definecolor{verbatimcolor}{rgb}{0.00,0.40,0.00}

\makeatletter

\renewcommand{\verbatim@font}
  {\ttfamily\footnotesize\selectfont}

\def\verbatim@processline{
  {\color{verbatimcolor}\the\verbatim@line}\par
}

\makeatother

%
% header
%

\rhead{}
\rfoot{\scriptsize{The kaneton microkernel project}}

\date{\scriptsize{\today}}


%
% header
%

\lhead{\scriptsize{The kaneton microkernel :: EPITA}}

%
% title
%

\title{The kaneton microkernel :: EPITA}

%
% authors
%

\author{\small{Julien Quintard \& Matthieu Bucchianeri}}

%
% document
%

\begin{document}

%
% title
%

\maketitle

%
% --------- text --------------------------------------------------------------
%

\begin{multicols}{2}

% XXX
\tableofcontents

%
% ---------- introduction -----------------------------------------------------
%

\section{Introduction}

kaneton is part of a larger project consisting in an educational
triology composed of:
  \textbf{k}\footnote{http://k.kaneton.org},
  \textbf{kaneton}\footnote{http://www.kaneton.org} and
  \textbf{kayou}\footnote{http://kayou.kaneton.org}; the
\textit{k* educational trilogy}.

Although these project remain educational, there exists an aside related
research project consisting in designing and implementing a novative operating
system.

Thus, people in charge of the teaching at \textit{EPITA} also contribute to
the development of the research project. Additionally, past teachers keep
contributing to the project either by developing or supervising the
new comers whilst some past students join the project.

%
% k
%

\subsection{k}

\textit{k} is an introductory course targeting low-level programming. The
project lasts several weeks and allows students to understand what is the
microprocessor's r\^ole in an operating system. Additionally, this course
allows the students to better unterstand the link existing between the
user-land applications and the operating system's kernel.

%
% kaneton
%

\subsection{kaneton}

\textit{kaneton} is probably the most important project of the trilogy. This
project, taught to second-year engineering students, lasts for about five
months. This project focuses on making students fully understand
internals of a micro-kernel based operating system.

%
% kayou
%

\subsection{kayou}

\textit{kayou} is a distributed operating system built over the kaneton
micro-kernel. The kayou's originality resides in its fully distributed
architecture. In such an architecture, all the machines of the network share
resources with each others. Shared resources include memory, processor,
storage, devices and so forth. kayou is not taught at EPITA, yet. However,
it would be a very good material for distributed systems courses.

%
% ---------- history ----------------------------------------------------------
%

\section{History}

During the kaneton history, the project evolved and courses were added
to the curriculum, making the whole educational project more interesting and
understandable to the students. Moreover, the educational project, which
was already targeting the \textit{EPITA}'s \textit{System, Network and
Security} major, was also used in other contexts, victim of its success
and of the very hard work achieved by kaneton people over the years.

%
% 2004
%

\subsection{2004}

The first year, an optional low-level programming introductory course named
\textbf{k} was proposed to the \textit{EPITA} Engineering School's first-year
students.

About fourteen hours courses were taught introducing the \textit{Intel 32-bit}
microprocessor's external architecture and low-level programming.

The students had to develop small, poor and messy device drivers for the
console and keyboard peripherals.

The course was a bit chaotic but this first shot was a success. Therefore, the
students majoring in \textit{System, Network and Security} asked the two
students \textit{Julien Quintard} and \textit{Jean-Pascal Billaud} for a
complete kernel project within their curriculum so that they can learn more
about operating systems internals.

%
% 2005
%

\subsection{2005}

The two, still students, \textit{Julien Quintard} and \textit{Jean-Pascal
Billaud} then started to work on a complete microkernel design the students
will have to implement. This was the premises of the \textbf{kaneton
microkernel educational project}. Additionally, two complete courses on kernels
design and \textit{Intel 32-bit Architecture} programming were prepared.

The project was composed of six steps, from the bootstrap, through the
kernel internals including memory management, task management etc.
to the servers with an \textit{IDE} device driver and finally a \textit{FAT}
file system.

The large majority of the students did not success in implementing a
complete scheduling system allowing the creation of user-land tasks. Indeed,
the best groups achieved in providing the management of kernel-land tasks only.
Therefore, the \textit{IDE} driver, \textit{FAT} file system etc. were
running inside the kernel, leading more to a monolithic kernel model.

Inspite of this, once again, the whole project was a success. However,
kaneton people noticed that the students took too much time doing unnecessary
time-demanding work like filling header files, dealing with versionning
problems, writing \textit{Make} files and \textit{Shell} scripts etc.

Moreover, the courses were too messy and the students had difficulties
to make the relation between the kaneton design and the microprocessor's
architecture implementation.

As a result, kaneton people decided to start implementing a kaneton microkernel
reference by their own in the \textit{C} language. This implementation will
then be used to compare the behaviour of students' implementation with
the reference. Moreover, this implementation led to the creation of a
new project: the \textbf{kaneton microkernel research project}.

%
% 2006
%

\subsection{2006}

While \textit{Jean-Pascal Billaud} leaved the project, people joined it
starting with \textit{EPITA} last-year students \textit{C\'edric Aubouy},
\textit{Renaud Lienhart} but also \textit{Fabien Le-Mentec} from
\textit{EPITECH} who knew all these people from the \textit{EPITA Computer
Systems Laboratory} where they were all working together a year before.
\textit{C\'edric Aubouy} and \textit{Renaud Lienhard} were in charge of the
kernel and \textit{Intel 32-bit Architecture} courses, respectively.
\textit{Julien Quintard} was still in charge of the kaneton educational
project taught to the students.

Moreover two \textit{EPITA} first-year students joined the \textit{EPITA
Computer Systems Laboratory}, \textit{Matthieu Bucchianeri} and \textit{Renaud
Voltz}. Indeed, from this date, the \textit{EPITA Computer Systems Laboratory}
was a strong partner of the kaneton microkernel project. These two students
were hired for contributing to the development of the kaneton research project.
Moreover, these students were supposed to teach and supervise the kaneton
educational project a year after.

\textit{Matthieu Bucchianeri} and \textit{Renaud Voltz} did an amazing work
on the kaneton research project implementation. Indeed, most of the
code related to the \textit{Intel 32-bit Architecture} comes from them. In
addition, the test suite as well as many tests were written by them.

This year, kaneton people decided to introduce a development environment,
based on the kaneton research reference implementation, including everything
necessary to set up a collaborative kernel development.

Whilst, previously, students had to write the entire microkernel and
servers from scratch, this year, students only had to write precise parts
of the microkernel including some set implementations, memory management,
task scheduling etc.

A few mistakes were made especially about the choice of parts the students
had to implement. Indeed, asking the students to implement set implementations
like linked-list, array etc. was a very bad idea. This year, the project
was not completed and students stopped the project before the messaging
system implementation.

A course was also added to the \textit{EPITA} \textit{System, Network
and Security} major's curriculum about microprocessors' internals. This
course was introduced and taught by \textit{Julien Quintard}.

In conclusion, the kaneton educational project was not a real success this year
and needed some modifications. For instance, the course about the \textit{Intel
32-bit Architecture} was too specific and hard to understand but also hard
to teach. Instead, kaneton people decided to replace it with a more general
course about kernel principles.

The kaneton research project implementation, in 2006,
  counted\footnote{Estimations realised with the software \textit{sloccount}.}
about \textit{7,000} lines for the \textit{core} and about \textit{2,000}
lines for the \textit{Intel 32-bit Architecture} implementation.

%
% 2007
%

\subsection{2007}

People affiliated with the \textit{EPITA Computer Systems Laboratory} joined
the project: \textit{Pierre Duteil} and \textit{Julian Pidancet}. Moreover,
students who implemented the kaneton educational project the previous year
decided to join the project: \textit{Enguerrand Raymond} and \textit{Mathieu
S\'elari\`es}, mainly working on the \textit{MIPS Architecture} portage amongst
other contributions.

This year, \textit{Matthieu Bucchianeri} and \textit{Renaud Voltz} were in
charge of the educational project by giving the kernel course as well
as supervising students' educational implementations.

As the kaneton research implementation was much more advanced as in 2006,
the students were given more code and then focused only on interesting and
system-related parts.

Students totally implemented the physical and virtual memory management, the
event and timer managers, the thread manager, the scheduler and the messaging
system. As for the previous years, the implementation was based on the
\textit{Intel 32-bit Architecture}.

This year, the whole kaneton educational project was also given to students
from the \textit{Realtime \& Embedded Systems} specialisation, for a total
of about $50$ students. The project was evaluated using a test suite,
developed for the kaneton research project, of about a hundred tests.

This year, the kaneton educational project was an amazing success as many
students completed a working microkernel, able to run tasks and to implement
servers running on top of the kaneton microkernel.

The kaneton research implementation has grown to \textit{9,000} lines of source
code for the \textit{core} and \textit{5,500} lines for the microprocessor's
architecture implementation on \textit{Intel 32-bit}. The kaneton research
implementation was able to start modules - as standalone binaries - in
user-land as well as to make them communicate through the kaneton messaging
system.

This year, \textit{Pierre Duteil} leaved the project prematurely. Since the
educational project needed two students for teaching kaneton to the second-year
students, \textit{Elie Bleton} joined the project.

This year was not as good as expected for the kaneton research project. As a
result, kaneton people decided to re-organise the curriculum of students
working on kaneton so that they can learn everything related to kaneton in a
very short period of time to be then able to work on their project without any
lack of knowledge.

%
% ---------- schedule ---------------------------------------------------------
%

\section{Schedule}

Until 2007, two students were recruited every year, in January, via the EPITA
Computer Systems Laboratory. These students were assigned a project and had
to work on it from February to the end of July.

The project assigned was, in the first years, directly related to the
development of the kaneton research project \textit{i.e} the reference
implementation. Through these projects, the recruited students discovered
and learned the concepts related to kaneton: micro-kernel design,
\textit{IA-32} microprocessor external architecture etc.

Unfortunately, this is no longer the case. Indeed, the development of the
kaneton reference reaches its end, the next steps being related to the kayou
distributed operating system. Thus, the new recruited students do not have
the opportunity to learn through their project.

Since students cannot learn about the kaneton design and \textit{IA-32}
programming, they will not be capable of teaching the kaneton courses and
project to the second-year students of EPITA.

As a result, these students have to be taught the fundamentals by more
experienced kaneton people. This reasoning led to the curriculum's schedule
described below.

The timetable below describes the teaching of the \textit{k* triology's}
projects at EPITA. Assistants are named \textit{Poussins} followed by the year
they have been recruited, \textit{Poussins'11} or \textit{Poussins$_{11}$} for
assistants recruited in 2011. For convenience and clarity, the whole year
$\gamma$ of the \textit{EPITA} curriculum is detailed next. As an example,
If $\gamma$ represents the year 2008, then, obviously,
\textit{Poussins$_{\gamma - 2}$} are assistants recruited in 2006.

%
% january \gamma
%

\subsection{January $\gamma$}

At the beginning of January, the \textit{EPITA Computer Systems Laboratory}
recruitment takes place. At the end of this period, a few students ---
ideally two or three --- would have been accepted to work on the kaneton and
kayou projects: the \textit{Poussins'$\gamma$}.

As explained above, these new comers need to be taught the fundamentals so
that they can work on their project with no lack of knowledge but also so
that they can teach and supervise the \textit{kaneton educational project}
in $\gamma + 1$. A one-week specific and very intensive period called
\textit{``piscine''} --- swimming pool in English --- is used to teach
the \textit{Poussins'$\gamma$} about everything related to kaneton.
Since these students have to acquire many notions in a very short period of
time, \textit{EPITA} is asked to arrange a special one-week curriculum for
these students.

This \textit{``piscine''} will be taugh by more experienced kaneton people,
especially the \textit{Poussins$_{\gamma - 2}$} possibly assisted by the
\textit{Poussins$_{\gamma - 1}$}.

At the end of this \textit{``piscine''}, \textit{Poussins} judged too weak
could be expelled. At the same time, projects are proposed to these students.
The \textit{Poussins} choose the project they will work on for the next
months. Every project, hence \textit{Poussin}, will be assigned a supervisor.
Note that supervisors can be \textit{EPITA Alumni} or people involved in
kaneton and/or kayou without any relation with \textit{EPITA}.

%
% february \gamma - june \gamma
%

\subsection{February $\gamma$ - June $\gamma$}

\textit{Poussins'$\gamma$} work on their assigned project in close relation
with their supervisor.

Depending on the progress status of the \textit{Poussins'$\gamma$} projects,
two students will be chosen for teaching and supervising the
\textit{k educational project} to the \textit{EPITA} first-year students.

If the \textit{Poussins'$\gamma$} supervisors consider that no student is
ready to carry out this task, then the courses will be taught by the
\textit{Poussins$_{\gamma - 1}$}, one-year more experienced.

%
% july \gamma
%

\subsection{July $\gamma$}

\textit{Poussins'$\gamma$} are asked to write and submit a report on the
work achieved through the five previous months. Additionally, every student
will have to give a talk, hence presenting their work to other students,
researchers, teachers etc.

This talk will be used by the head of \textit{EPITA Computer Systems
Laboratory} for assigning a student mark which will be used by the
\textit{EPITA Board of Studies}.

Once evaluated, two \textit{Poussins'$\gamma$} are picked for carrying out,
in February, the teaching and supervising of the \textit{kaneton educational
project} to the second-year students major in both \textit{System, Network and
Security}; and \textit{Embedded and Real-Time Systems}.

%
% august \gamma - january \gamma + 1
%

\subsection{August $\gamma$ - January $\gamma + 1$}

Whilst interns, \textit{Poussins'$\gamma$} are due to prepare the teaching
starting in February. This preparation includes the source code, tests,
lectures, subjects, documentation as well as any administrative task: budget,
timetable and so forth.

%
% february \gamma + 1 - june \gamma + 1
%

\subsection{February $\gamma + 1$ - June $\gamma + 1$}

The selected \textit{Poussins'$\gamma$} teach and supervise the second-year
students through the \textit{kaneton educational project}.

%
% july \gamma + 1 - january \gamma + 2
%

\subsection{July $\gamma + 1$ - January $\gamma + 2$}

The \textit{Poussins'$\gamma$}, before leaving \textit{EPITA} for their
last-year internship, can take some time to contribute to the \textit{kaneton}
and/or \textit{kayou} research projects.

%
% february \gamma + 2 - \infinity
%

\subsection{February $\gamma + 2$ - $\infty$}

The \textit{Poussins'$\gamma$} are more than welcome to keep contributing
to any project of the \textit{k* educational trilogy} or to one of the related
research projects.

%
% ---------- k* educational triology ------------------------------------------
%

\section{k* Educational Trilogy}

%
% k
%

\subsection{k}

k is a short but complete introduction to low-level programming and
operating system basics.

Students - in groups of four - develop a tiny monolithic kernel
intended to run very simple applications in a mono-task environment.

While dealing with many fundamental aspects of operating systems like
device drivers and system calls, the project is funny and visual. The
programs to run are tiny 2D-arcade games with graphic rendering and
sound.

% project

\subsubsection{Project}

The project is divided into several parts:

\begin{itemize}
  \item
    \textbf{pm} manager that takes care of memory management;
  \item
    \textbf{event} and \textbf{syscall} managers handling external
    events and user-land to kernel-space transitions;
  \item
    \textbf{cons}, \textbf{vga} and \textbf{speaker} managers for
    multimedia operations;
  \item
    \textbf{kbd} and \textbf{mouse} managers to provide user inputs;
  \item
    \textbf{time} manager to measure time intervals;
  \item
    \textbf{fs} manager to manage the root filesystem.
\end{itemize}

The project lasts two months and is concluded by an evaluation of the
final projects with each groups.

% courses

\subsubsection{Courses}

The project comes with a single course.

% XXX[cette section devrait etre completee avec 1) les points abordes 2)
%     un detail de chaque session. voir MIA et KER ci dessous pour un modele]

\textit{Name}: \textbf{K}

\textit{Hours}: 14 hours divided into 6 sessions

\begin{itemize}
  \item
    Introduction to memory management, especially the segmentation mechanism;
  \item
    Presentation of the \textit{IA-32} calling convention;
  \item
    Event handling, including hardware interrupts;
  \item
    Privilege modes and system calls;
  \item
    Device driver implementation and filesystems;
  \item
    Binary formats and ELF loading.
\end{itemize}

%
% kaneton
%

\subsection{kaneton}

The kaneton microkernel project aims at leading students to understand
what really is an operating system and more precisely a kernel.

Every group, composed of two students, has to develop parts of the kaneton
microkernel.

To simplify the project, each group receives a development environment
containing everything necessary to start the project including a
working microkernel base, documentation, assignments etc.

% project

\subsubsection{Project}

The whole kaneton project is subdivided into stages:

\begin{itemize}
  \item
    \textbf{k0}: this project will be entirely developed during a
    special practical session of about 6 hours and under the teachers'
    supervisation. This project consists in an introduction to the
    IA-32 external architecture: basic memory management model,
    IA-32 assembly language, development tools etc; (1 day)
  \item
    \textbf{k1}: this project will lead the students to an higher abstraction
    level. The goal of the project is to develop: the segment manager to be
    able to manage physical memory; and the architecture-dependent part of the
    region manager to be able to control the virtual memory; (2 weeks)
  \item
    \textbf{k2}: this project introduces the interrupts management and
    the driver development. Indeed, the student will have to set up
    everything necessary to handle interrupts. Then, some drivers will
    be developed including the keyboard and timer drivers as well as a
    demonstration shell; (2 weeks)
  \item
    \textbf{k3}: the students have to provide everything necessary
    to handle execution contexts. Indeed, the student will have to
    develop the scheduler but also to implement the context switchs from
    the microprocessor point of view; (3 weeks)
  \item
    \textbf{k4}: this project introduces the communication into the
    microkernel. Indeed, the students will have to provide some
    communication primitives including sychronous and asynchronous
    functions as well as blocking and non-blocking functions. (2 weeks)
\end{itemize}

% courses

\subsubsection{Courses}

The project comes with two fundamental courses:

\begin{itemize}
  \item
    \textbf{MIPS Architecture}: this course is dedicated to the MIPS internal
    architecture although this course actually explains modern microprocessors
    concepts: pipelining, optimizations, caches etc;
  \item
    \textbf{Kernel}: this course explores the kernels internals including
    microkernel architectures, memory management, processes and threads,
    synchronisation, communications, and complete case studies.
\end{itemize}

% _mips architecture_

\textit{Name}: \textbf{MIA}

\textit{Hours}: 30 hours divided into 10 sessions

This course is intended to give students a complete view of how a RISC
microprocessor works and what choices the designers made to build the
future modern microprocessors.

Below are listed the concepts studied in this course:

\begin{itemize}
  \item
    External architecture: the instructions set, the memory management model
    etc;
  \item
    Pipeline: its inherent problems and limitations. Advanced pipelines will
    also be studied;
  \item
    Compiler optimisations: direct assembly source code optimisations:
    rescheduling, software pipelining etc;
  \item
    Memory: bus used on some experimental processors and the different cache
    management techniques.
\end{itemize}

The list below details the different sessions for the MIPS architecture
course:

\begin{enumerate}
  \item
    Course presentation, MIPS history, introduction to MIPS external
    architecture: registers and instructions set. Introduction to
    instruction formats and to inherent limitations; (3 hours)
  \item
    End of instruction formats introduction, explanations of architecture
    designers choices: MIXs and Amdhal Rule. Some exercises to practice.
    Introduction to MIPS architecture addressing and to MIPS pipeline;
    (3 hours)
  \item
    Introduction to MIPS spirit and to its pipeline. Moore Law and
    pipeline rules. Introduction to different pipeline representations:
    simplified and detailed. Some exercises to practice these representations;
    (3 hours)
  \item
    We will study in this session the next instruction address computation
    problem and the delay slot. Then, study of branch instruction and their
    problems and limitations. Some exercises to practice; (3 hours)
  \item
    Complete study of instruction dependencies. Some exercise to practice;
    (3 hours)
  \item
    Partial exam correction. Introduction to optimisations; (3 hours)
  \item
    Study of different optimisations used by compilers. Study of
    advanced pipelines: superpipeline, superscalar pipelines etc; (3 hours)
  \item
    Introduction to the Pi-Bus. Then complete study of memory including
    cache management; (3 hours)
  \item
    End of memory course. Some exercises to practice; (3 hours)
  \item
    Revisions for the final exam. (3 hours)
\end{enumerate}

% _kernel_

\textit{Name}: \textbf{KER}

\textit{Hours}: 26 hours divided into 9 sessions

This course is intended to give the students a general overview of
kernel internals and its problematics.

Below are listed the studied notions:

\begin{itemize}
  \item
    Understand how kernel developement evolves and why, receive a
    minimal general culture on kernel internals and history;
  \item
    Resources --- memory, execution time --- sharing and management with
    increasing constraints and needs;
  \item
    The microprocessor's facilities are detailed to get the students
    to fully understand the mechanisms. Examples are taken from
    many different architectures;
  \item
    The common kernel design issues will be totally understood by the
    students and they will be able to build their own design;
  \item
    The course aims at giving the students a precise idea of how it
    works inside, thus giving them the ability to fully understand
    the kaneton project.
\end{itemize}

The list below details the different sessions for the kernel course.

\begin{enumerate}
  \item
    Recalls and prerequisites: presentation of the diffent kind of kernels
    and the role in an operating system. Recalls of low-level basics; (3 hours)
  \item
    Bootstrap: details about processor bootup and bootloading; (2 hours)
  \item
    Memory management - part one: fundamental concepts of memory management,
    overview of memory management unit and address translation though paging;
    (3 hours)
  \item
    Memory management - part two: high-level mechanisms like page swapping,
    copy-on-write, etc. Presentation of most known allocation algorithms;
    (3 hours)
  \item
    Events and I/O - part one: interrupt mechanism, interrupt vector and
    processor running modes; (3 hours)
  \item
    Events and I/O - part two: hardware interrupts, interrupt controllers,
    port-mapped input/output, memory-mapped input/output; (3 hours)
  \item
    Task and scheduling - part one: recalls on multiprogramming, execution
    context, context switching; (3 hours)
  \item
    Task and scheduling - part two: process control block on known operating
    systems, task life cycle, scheduling automaton, scheduling algorithms;
    (3 hours)
  \item
    Inter-Process Communication: \textit{IPC} in a microkernel, message
    passing, classical \textit{IPC}, locking. (3 hours)
\end{enumerate}

%
% kayou
%

\subsection{kayou}

kayou is not taught at \textit{EPITA} yet.

%
% ---------- project suggestions ----------------------------------------------
%

\section{Project Suggestions}

This section contains some project suggestions for the new \textit{Poussins}.

%
% security
%

\subsection{Security}

Security is a very important topic in operating systems, even more in
distributed systems. Indeed, in such a system, resources are completely
shared amongst the machines of the system. It is therefore crucial to protect
these resources from being accessed by unauthorised entities.

This project consists in designing and implementing a complete security
model composed of some in-kernel primitives as well as at the operating
system level.

For now, the kaneton microkernel provides the concept of \textit{capabilities}
which seems to perfectly fit the distributed systems' requirements. The student
will have to confirm that this concept is suitable for kayou's needs. The
student will also have to complete the current implementation in order to move
to the study of the security in the operating system, even its in distributed
form.

A possibly suitable access control scheme could be the \textit{RBAC} -
Role-Based Access Control. The student will have, once again, to confirm this
possibility via a preliminary study and then implement the services
necessary to this scheme so that applications can use it.

The student will have to tackle the notion of multi-users at the operating
system level. Indeed, although this notion is integrated at the kernel
level in every \textit{UNIX} systems, kaneton designers decided to handle this
notion outside the microkernel.

The student will then have to study this concept and implement the necessary
services. Besides, the student could also investigate further in order
to develop a security service based on the concept of user. \textit{RBAC}
could, once again, turn out to be a good choice.

%
% posix
%

\subsection{POSIX}

kaneton provides its own interface which has nothing to do with interfaces
of others \textit{UNIX} like \textit{Linux}, \textit{BSD} or others etc.

\textit{POSIX} is an interface which was introduced in order to normalise the
\textit{C Library} layer. Most applications --- at least in the \textit{UNIX}
world --- rely on this software layer.

If we want our operating system to benefit from the --- already existing and
future --- software without having to re-develop everything, we need to
develop a software layer complying to the \textit{POSIX} standard. In our case,
the \textit{POSIX} functionalities must be translated into either
functionalities provided either by kaneton or one or more kayou servers.

The student will have to study very precisely the \textit{POSIX} standard
as well as the kaneton interface in order to develop the emulation library:
\textit{libposix}.

In order to test the library, the student will be in charge of porting some
services and applications. Since \textit{UNIX} highly relies on the concept
of \textit{files}, it is very likely that the student will have to port
or develop a hard-disk driver as well as a basic file system.

The student could take advantage of the different work existing on the
Internet, especially \textit{LSE/OS} which was designed and implemented
at the EPITA Computer Systems Laboratory.

%
% amd64
%

\subsection{AMD-64}

This project consists in porting kaneton to a new microprocessor architecture,
in this case, \textit{AMD-64}.

This architecture is very interesting for two reasons. First, \textit{AMD}
architectures have many commonalities with \textit{IA-32} architectures. Since,
kaneton already supports \textit{IA-32}, it would be easy to support this new
architecture. Secondly, this architecture runs in 64-bit and its implementation
could reveal design mistakes and/or implementation bugs in the generic
part of the microkernel, commonly called \textit{core}.

The student will have to study a complete new architecture by reading the
microprocessor specifications. Besides, the student will have to report
any problem encountered in his project, especially about the microkernel's
design. For example, types handling in the \textit{core} is very likely to
be problematic.

The student will have to modify several parts in the kernel but also
make suggestions about design improvements based on his observations.

Besides, the student will have to complete the microkernel implementation
in various ways: functionalities not yet implemented, bugs correction, error
handling improvement, etc.

%
% mips-64
%

\subsection{MIPS-64}

This project simply consists in contributing to the portage of kaneton
on the \textit{MIPS-64} architecture.

The remarks about the \textit{AMD-64} project apply here as well.

%
% dht
%

\subsection{DHT}

This project consists in developing a \textit{DHT} --- Distributed Hash Table.
This abstraction is crucial for distributed applications as these applications
need to store data and make them accessible to the other nodes of the
distributed system.

The \textit{DHT} is built over what is called a \textit{KBR} --- Key-Based
Routing. The r\^ole of the \textit{KBR} is to, given a key, locate the
node which is in charge for storing a piece of data.

This \textit{DHT} will be heavily used in kayou, especially for the distributed
file system, but could also be used by many distributed applications.

There does not exist any good implementation of such \textit{DHT}s. Indeed,
the well-known implementations were written in Java and are therefore
useless in a true system where performances matter.

The student will have to briefly study this research field through the
well-known papers and projects: Chord, \textit{CAN}, Pastry etc. and then
to design and finally implement the complete \textit{DHT}.

This project is clearly research-oriented since the concepts directly come
from research. However, many firms are highly interested in such a
performant and secure implementation in order to build large-scale storage
systems. Note that the student will have the opportunity to write a paper and
possibly present the work achieved on this project.

The student will probably contribute to the development of a peer-to-peer
file system in order to test the \textit{DHT} but also to contribute to
another kayou project.

\end{multicols}

\end{document}

%
% ---------- XXX --------------------------------------------------------------
%

- planifier les cours/tp etc. en fonction des terms, des tranches etc.
  pour que les exams tombent bien! en deduire la longeur des tranches.
- prevoir le budget.
- reunion avec Dudu avec toutes les informations dans des documents Excel:
  heures cours, heures perm, nombre exams etc. + heures de travail par etudiant
  par tranche et faire accepter le budget (negocier! mais 1. il accepte le
  budget en fonction de la confiance qu'il accorde au profs. il faut donc
  lui faire un dossier niquel pour montrer que vous etes serieux et lui
  faire oublier ce qu'il a pu penser de mal sur vous 2. puisque le projet
  est deja quasiment tout pret: cours, exams, tarball, intra etc. possible
  qu'il veuille baisser le budget. a vous de trouver des arguments si vous
  pensez que c'est abuse!).
- prendre en main les cours: les relire, se les appropriez, les retoucher
  et/ou completer. ceci inclut les cours noyau mais aussi les cours de remise
  a niveau etc.
- prendre en main les presentations de projets.
- prendre en main kaneton: design, systeme de portabilite, environnement etc.
  pour etre capable de repondre aux questions des etudiants.
  genre sur le wiki ou dans un document, ecrire quelques lignes sur ``Comment
  mettre en place votre environnement kaneton en 3 minutes'': var d'env,
  creer un profile user (optionnel), make initialize, make build, make install
  avec peut etre les pitfalls.
- prendre en main l'intra LSE, le reinitialise, savoir comment l'editer,
  ajouter des users (les novueaux etudiants), uploader les resultats de
  moulette, recuperer tarball uploadees etc.
- prendre en main la moulette pour faire tourner les tests et obtenir des
  traces utilisables.

---

la suite doit aller sur le wiki Teachers

% preparation

\subsubsection{Preparation}

This section lists --- in a more or less chronologic manner --- the tasks that
new kaneton teachers must undertake before February.

\begin{itemize}
  \item
    Plan the sub-projects the students will have to complete probably from
    \textit{k0} to \textit{k4};
  \item
    According to the previous planning, update the kaneton research project's
    source code so that tags are well placed.

    \-

    The student tarball could then be generated via the kaneton
    \textit{export} tool;
  \item
    Plan the lectures, practical sessions, exams and surgerie hours.

    \-

    Especially, a general presentation about the upcoming project is held
    at the beginning of the term, in February. Additionally, few hours courses
    are given to the students focusing on the project's requirements:
    \textit{C} language, \textit{Assembler} language, \textit{inline Assembler}
    etc. Finally, the \textit{k0} project is generally performed as a practical
    session;
  \item
    XXX
\end{itemize}

